\documentclass[10pt]{article}
\usepackage[margin=1in]{geometry}
\usepackage{booktabs}
\usepackage{longtable}
\usepackage{array}
\usepackage{fontspec}
\XeTeXlinebreaklocale "zh"
\XeTeXlinebreakskip = 0pt plus 1pt
\IfFontExistsTF{Noto Serif CJK SC}{\setmainfont{Noto Serif CJK SC}}{\setmainfont{Microsoft YaHei}}
\title{标题}
\author{path-scientist-agent}
\date{}
\begin{document}
\maketitle

PathMNIST数据集上SmallResNet的优化与评估：从验证集到测试集的泛化差距分析
\section{摘要}
本文介绍了在PathMNIST数据集上对SmallResNet模型进行的训练、调优和单次测试评估。研究严格遵循了仅使用训练/验证集进行模型选择和超参数调优的原则，并在候选模型冻结后对测试集进行了且仅进行了一次评估。在训练和调优阶段，我们比较了多种变体（基线、数据增强、多尺度、优化及组合），其中“优化”变体在验证集上取得了0.997589的最佳Macro-F1分数。然而，在最终冻结后的单次测试集评估中，该模型的Macro-F1均值下降至0.834035，验证集与测试集之间存在0.158425的显著差距。具体而言，种子27在测试集上表现出不稳定性，显著降低了整体均值并增加了方差。本文不提出任何临床使用或诊断声明。
\section{引言}
PathMNIST是一个广泛用于医学图像分类的数据集。在机器学习研究中，确保模型评估的严谨性至关重要，特别是在避免测试集信息泄露方面。本研究旨在探讨SmallResNet模型在PathMNIST数据集上的性能表现，重点关注超参数优化策略以及模型从验证集到测试集的泛化能力。我们严格遵循了在模型冻结后仅对测试集进行一次评估的原则，以提供无偏见的泛化性能估计。本研究不涉及任何临床应用或诊断声明。
\section{相关工作}
在深度学习领域，残差网络已被证明在图像分类任务中具有卓越的性能。针对医学图像数据集如PathMNIST，研究人员通常采用数据增强、多尺度训练以及学习率调度等技术来提升模型性能。然而，模型在验证集上的优异表现并不总是能转化为测试集上的同等性能，泛化差距是一个常见问题。本文在此基础上，系统评估了SmallResNet在不同配置下的表现，并报告了冻结候选模型在测试集上的单次评估结果。
\section{材料与方法}
\textbf{数据集}: PathMNIST (SHA-256: \texttt{1e7fc200dd5aac79f39f0da26178be801ffac8e6d59c8533e51acebae953745a})。
\textbf{模型}: SmallResNet。
\textbf{变体}: 优化。
\textbf{超参数}: 学习率为0.001，权重衰减为1e-05，采用OneCycle学习率策略，标签平滑系数为0.1，未使用数据增强和多尺度技术。
\textbf{随机种子}: 7, 17, 27。
\textbf{调优过程}: 在冻结前完成了六种网格设置，仅使用训练/验证集进行调优。最佳学习率和权重衰减分别为0.001和1e-05，最佳验证Macro-F1为0.997589。
\textbf{评估策略}: 候选模型冻结后，在测试集上进行了且仅进行了一次评估。测试结果未用于调优、模型选择或重新训练。
\section{实验}
在实验阶段，我们首先在训练/验证集上进行了多组对比实验，包括基线、数据增强、优化、多尺度及其组合与消融实验。每种配置均运行3个种子。完成调优后，确定了“优化”变体为最终冻结候选模型。随后，在测试集上对该候选模型进行了单次评估。运行环境为PyTorch 2.13.0+cu132，CUDA 13.2，GPU为NVIDIA GeForce RTX 5070 Ti Laptop GPU。
\section{结果}
在训练/验证阶段，“优化”变体取得了最佳表现，Macro-F1均值为0.992459，标准差为0.008029。其他变体的验证集结果如下：基线(0.956763)，数据增强(0.971311)，多尺度(0.978622)，组合(0.968218)。消融实验中，combined\_no\_augmentation为0.930088，combined\_no\_multiscale为0.949365。
在最终单次测试集评估中，测试Macro-F1均值为0.834035，标准差为0.075947。验证集到测试集的Macro-F1差距为0.158425。
各种子测试结果如下：
\begin{itemize}
\item 种子7: 最佳epoch为39，准确率为0.906128，Macro-F1为0.875924，损失为0.397548。
\item 种子17: 最佳epoch为39，准确率为0.906267，Macro-F1为0.879813，损失为0.412215。
\item 种子27: 最佳epoch为18，准确率为0.779666，Macro-F1为0.746367，损失为0.722006。
\end{itemize}
在聚合混淆矩阵中，召回率最低的三个类别分别是类别7 (0.358)、类别0 (0.724)和类别6 (0.858)。
\section{讨论}
实验结果表明，仅使用优化策略的候选模型在从验证集迁移到测试集时存在显著的泛化差距（0.158425）。尽管在验证集上取得了接近完美的Macro-F1（0.997589），但测试集上的表现大幅下降。特别是种子27在测试集上表现出明显的不稳定性，其Macro-F1仅为0.746367，这显著拉低了整体的均值并增加了方差。聚合混淆矩阵的分析显示，类别7、0和6的召回率最低，表明模型在这些特定类别上存在混淆。这些观察结果仅作为最终报告，未触发进一步的调优或模型修改。
\section{局限性}
本研究的局限性在于测试集评估仅进行了一次，且仅使用了三个随机种子。种子27的异常不稳定性对整体结果的均值和方差产生了较大影响。此外，模型未采用数据增强和多尺度技术，这可能限制了其在某些困难类别上的泛化能力。本文不提出任何临床使用或诊断声明。
\section{结论}
本研究严格遵循了仅在训练/验证集上进行调优、在模型冻结后对测试集进行单次评估的实验规范。SmallResNet的“优化”变体在PathMNIST验证集上表现优异，但在测试集上表现出显著的泛化差距。种子27的不稳定性以及特定类别（7、0、6）的低召回率是导致测试性能下降的主要因素。研究结果强调了验证集与测试集之间泛化差距的重要性。
\section{可复现性声明}
为确保本研究的完全可复现性，以下信息和工件已记录在案：
\begin{itemize}
\item 数据集SHA-256: \texttt{1e7fc200dd5aac79f39f0da26178be801ffac8e6d59c8533e51acebae953745a}
\item 运行环境: PyTorch \texttt{2.13.0+cu132}, CUDA \texttt{13.2}, NVIDIA GeForce RTX 5070 Ti Laptop GPU
\item 相关工件路径:
\item \texttt{runs\textbackslash{}pathmnist-m4\textbackslash{}tuning\textbackslash{}result.json}
\item \texttt{runs\textbackslash{}pathmnist-m4\textbackslash{}final\textbackslash{}optimization\textbackslash{}seed\_7\textbackslash{}checkpoint.pt}
\item \texttt{runs\textbackslash{}pathmnist-m4\textbackslash{}final\textbackslash{}optimization\textbackslash{}seed\_17\textbackslash{}checkpoint.pt}
\item \texttt{runs\textbackslash{}pathmnist-m4\textbackslash{}final\textbackslash{}optimization\textbackslash{}seed\_27\textbackslash{}checkpoint.pt}
\item \texttt{runs\textbackslash{}pathmnist-m4\textbackslash{}test\_evaluation.json}
\item 测试集评估在候选模型冻结后仅进行了一次，未用于任何后续调优或模型选择。
\end{itemize}

\end{document}
